\subsection{Restricciones de Diseño}
% Contenido pendiente de redacción.
El diseño de la aplicación deberá respetar las siguientes restricciones, cuyo detalle se documentará en la arquitectura y los modelos UML del proyecto~\cite{pmp2026}:

\begin{itemize}
  \item \textbf{Tipo de aplicación y alcance.} El producto deberá ser una aplicación de escritorio independiente para Windows~11. En esta primera versión solo cubrirá el sector de Aguas Claras, de modo que su uso en otros sectores dependerá de ampliar la biblioteca de escenarios.
  \item \textbf{Simulador.} Los escenarios de la biblioteca deberán generarse con FDS~6.11.1~\cite{fdsguide} y su modelo de vegetación, a partir de las capas de terreno y vegetación preparadas con QGIS~4.2 y qgis2fds~2.0.1. Las simulaciones se ejecutarán en HPC-ZINE durante la construcción del producto, y la aplicación solo leerá sus resultados.
  \item \textbf{Lenguaje de programación y bibliotecas.} La aplicación deberá desarrollarse en Python~3.13, con PySide6~6.12 para la interfaz gráfica, fdsreader~1.14 para leer las salidas de FDS y Shapely~2.2 y pyproj~3.8 para las operaciones geométricas y la transformación de coordenadas~\cite{pmp2026}.
  \item \textbf{Restricciones por el lenguaje.} La aplicación deberá empaquetarse con PyInstaller~6.22 como un ejecutable para Windows, que incluye el intérprete de Python y sus bibliotecas; por ello el ejecutable será más pesado que uno nativo y solo funcionará en el sistema operativo para el que se construya~\cite{pmp2026}.
  \item \textbf{Arquitectura.} La aplicación no seguirá una arquitectura cliente-servidor: se ejecutará localmente, organizada en los cuatro módulos de la sección~2.7, y consultará la biblioteca de escenarios como un repositorio local. Su única comunicación externa será la consulta a Open-Meteo.
  \item \textbf{Modelo de la base de datos.} La biblioteca deberá usar una base de datos relacional en SQLite, construida a partir del modelo del dominio de la sección~2.5; su indexación y su búsqueda se especifican en la sección~3.6.
  \item \textbf{Herramientas CASE.} Los modelos de análisis y diseño deberán elaborarse en UML~2.5.1 con PlantUML, e incluirán como mínimo los diagramas de casos de uso, componentes y despliegue~\cite{pmp2026}.
\end{itemize}